% mage-auto-heading: cloudflare
% mage-auto-heading: spotify
% mage-auto-heading: shopify
% mage-auto-heading: uber
% mage-auto-heading: amazon
% mage-auto-heading: gitlab
% mage-auto-heading: what-the-factories-share-and-where-they-differ
% Authoritative LaTeX manuscript; Markdown is a generated compatibility projection.

% <!-- part-title: Agentic Software Factories -->

% <!-- chapter-title: Other Agentic Software Factories -->

\chapter{Other Agentic Software Factories}

% <!-- point: one-system-one-engineer-one-problem-so-how-do-others-do-it | Public industrial accounts broaden the one-engineer case by showing different arrangements of knowledge, production, evidence, and authority; the comparison uses the four factory questions. -->

DocAble showed how one engineer assembled a factory around coding agents. What changes when agents work inside an established engineering organization, or when a team adopts machinery built elsewhere? The four questions from \magevalue{sec:the-software-factory} guide the comparison: what the factory produces, what people do, how production is organized, and what tolerances the resulting software must satisfy. The accounts below illuminate different parts of those questions. They need not supply the same answer, or document every part equally well.

% <!-- point: the-public-sources-are-interested-accounts-read-conservatively | The cases are interpretive reconstructions from interested public sources, not independent evaluations or claims of MAGE adoption; mechanisms, measurements, and causal interpretations have different evidential status. -->

These are my interpretive reconstructions of public reports, not claims that the companies adopted MAGE. Engineering blogs, talks, and product documentation describe mechanisms and operating experience, but also serve promotional purposes. I use MAGE to interpret the arrangements they describe; the reports do not independently validate the method. Reported measurements retain their scope, and missing information remains missing rather than proving that a mechanism does not exist.

\section{Cloudflare}\label{heading:cloudflare}

% <!-- point: cloudflare-turns-standards-prose-into-statements-with-stable-identity | Cloudflare represents institutional obligations as individually identified statements that can be tracked through revision and applied at several stages of engineering work. -->

Cloudflare shows how an organization\textquotesingle s policy can become a control over work performed across many teams. Its standards corpus had grown beyond what individual engineers could readily remember. An extraction agent turns each SHOULD or MUST into a structured statement with a stable identifier, so the obligation can be tracked when its containing document changes.\citep{reimann2026codex} That identity supports compliance records and explicit waivers. Standards then meet work at several stages: specification review before implementation, linters during editing, semantic review in CI, and incident review after production. The representation connects a human-authored obligation to machinery that can find and apply it; merely storing the original policy documents would leave more of that connection to each reader.

% <!-- point: cloudflare-promotes-obligations-into-blocking-authority-and-instruments-one-side-only | Standards approval and enforcement promotion are separate; enforced MUST statements can drive merge decisions, with a recorded human override, while published activity counts do not establish detection performance. -->

An approved statement does not automatically block work. Any employee may propose a standard, and a named domain owner approves it. A separate promotion moves it from advisory to enforced; the cited account does not identify who performs that step. For enforced standards, a coordinator agent orchestrates up to seven specialist reviewers, evaluates severity, and can approve a merge request, revoke an earlier approval, or block the merge. Only enforced MUST statements can block, and the rubric favors approval so that a lone warning does not.\citep{skidmore2026review} This separates the decision to adopt an obligation from the decision to let an automated reviewer enforce it.

Human reviewers retain an exception path: commenting ``break glass'' forces approval and records the override. Cloudflare reports 288 such overrides, 0.6 percent of merge requests, from March 10 to April 9, 2026.\citep{skidmore2026review} That is useful evidence about the operation of the control. It does not estimate violations the reviewers missed or defects they prevented. Review runs, findings, and blocked merges likewise measure activity. The account makes the allocation of authority visible even where the effectiveness of the checking remains incompletely measured.

% <!-- point: cloudflare-converts-failures-into-structure-but-never-measures-the-loop | Named incidents and operating costs changed rules, shared deployment machinery, representations, and check placement; the reports do not measure the resulting reduction in failures. -->

Operating experience also changes what later work inherits. The November and December 2025 outages led to rules against calling \texttt{.unwrap()} outside tests and against assuming upstream dependencies were in an expected state. Dangerous configuration patterns can be enrolled in Snapstone\textquotesingle s shared deployment machinery. The policy representation changed from Markdown to JSON when agents needed more reliable filtering.\citep{hartman2026failsmall} Semantic-review latency also encouraged moving suitable checks into deterministic linters earlier in the workflow.\citep{reimann2026codex} These are concrete instances of governance conversion: a failure or recurring cost leads to a reusable change in the environment. The sources describe those changes but do not measure their effect on later incidents.

\section{Spotify}\label{heading:spotify}

% <!-- point: spotify-holds-the-line-fixed-and-swaps-the-station-inside-it | Spotify introduced an agent into a pre-existing fleet-change process, allowing examination of reused production structures without claiming causal isolation. -->

Spotify introduced agents into a production process that already knew how to change a large software estate. Fleet Management selected repositories, ran transformations, opened and merged pull requests, and monitored the resulting deployments, alerting the change author when production failed.\citep{brown2023fleet} In 2025, an agent replaced the transformation script at one station while the surrounding targeting, review, and merge infrastructure initially remained the same. The agent addressed changes that deterministic scripts had struggled to cover, including the remaining 30 percent of one fleet migration.\citep{charas2025honk} This lets us examine which production structures could be reused and which later required revision.

% <!-- point: spotify-installs-and-then-retires-a-control-as-the-fabricator-improves | Bounded tools and mandatory checks govern proposed changes; an added judge and its reported later removal show controls changing with observed failures and false vetoes. -->

The late-2025 configuration bounded both the agent\textquotesingle s actions and its route to a pull request. It exposed selected Git commands that excluded pushing, shell commands from an allowlist, and a verification tool. Users supplied relevant context in prompts; code-search and documentation tools were withheld.\citep{charas2025honk} The verification tool selected checks from repository contents, and failure prevented a pull request from being opened. The agent could choose how to implement a change within those bounds, but its own assessment could not waive the required checks. After unwanted changes such as disabled tests, Spotify added an LLM judge that compared the diff with the original request. In December 2025, the company reported vetoes in about a quarter of thousands of sessions, with agents correcting their work after about half of those vetoes.\citep{charas2025feedback}

A March 2026 report of Spotify\textquotesingle s QCon talk describes a further separation between the agent runtime and verification runtime, with pull requests created only after CI checks. It also reports that the judge sometimes rejected valid changes and was later removed as foundation models improved and verification instructions in prompts proved sufficient for its role.\citep{curtis2026spotifyqcon} This is a secondary account of the talk, and it does not imply that the required checks or review disappeared. The sequence shows why maintaining a control includes reassessing both the failures it catches and the valid work it rejects. An engineering response to one generation of agent behavior need not remain economical unchanged.

% <!-- point: spotify-splits-admission-by-fabricator-and-names-verification-as-the-constraint | Deterministic fleet changes and agent-authored changes follow different admission arrangements; reported operating experience exposes verification capacity as a constraint. -->

The inherited process also supports different admission decisions for different producers. Spotify reports more than 2.5 million automated maintenance pull requests, most automatically merged; that cumulative figure chiefly describes the deterministic fleet machinery, not agent-authored work.\citep{spotify2026constraint} The QCon account says agent-authored changes still required human review.\citep{curtis2026spotifyqcon} People define migration scope and maintain prompts as tested, version-controlled artifacts. They can also reserve particular decisions: in one migration, the agent was instructed to leave fields unchanged and annotate them whenever a judgment call was needed.\citep{charas2025honk, edwardsjoseph2026migrations} Delegation therefore varies within the production process rather than being one decision to automate the whole line.

% <!-- point: spotify-measures-the-health-of-the-engineering-system-not-just-agent-activity | Rework, complexity, pull-request size, and incident observations provide a mixed account of system health and verification capacity rather than an activity-only success claim. -->

Spotify\textquotesingle s later measurements make the limits of that process visible. Its refined rework measure did not rise during the AI ramp-up, while code complexity and pull-request size did; the company kept its thresholds unchanged. Incident retrospectives found no material direct contribution from AI-authored code, yet change volume had grown faster than some verification controls could adapt. One automated dependency upgrade passed its checks and still failed in production.\citep{singer2026quality} These observations concern different properties and should not be collapsed into a single quality score. Producing more changes has exposed a capacity problem as well as a check-coverage problem.

\section{Shopify}\label{heading:shopify}

% <!-- point: shopify-makes-every-session-observable-organizational-material | Public River sessions let colleagues inspect work and failures, creating opportunities to record knowledge for later sessions without establishing that every failure is watched. -->

Shopify\textquotesingle s case concerns how work performed in one session becomes knowledge available to later sessions. River, its internal agent, operates in Slack and refuses direct messages. An employee starts work by mentioning it in a channel, leaving a public transcript that colleagues can inspect.\citep{moreno2026river} This makes the work and its difficulties observable beyond the person who initiated it. It does not ensure that every failure attracts attention, but it provides a shared record from which people can identify missing knowledge instead of depending entirely on a later report.

% <!-- point: shopify-mines-sessions-by-hand-and-measured-the-effect | People record missing knowledge in shared files; the reported merge-rate improvement occurred with the foundation model held fixed but does not isolate human learning from environmental changes. -->

People use that record to change what later agents receive. They watch River work, identify what it should have known, and add a skill, revise an \texttt{AGENTS.md} file, or record another reusable instruction in the monorepo.\citep{moreno2026river} These artifacts sit alongside code, conventions, intent documents, and runbooks. River\textquotesingle s reported merge rate rose from 36 to 77 percent over two months while its foundation model stayed the same.\citep{lutke2026shopfloor} The observation does not separate human learning from changes to the environment. It does show an improvement without a foundation-model change, in an arrangement where people deliberately preserve lessons for later work. Unlike DocAble\textquotesingle s one-engineer account, this mechanism distributes observation and knowledge maintenance across an organization.

% <!-- point: shopifys-instruction-corpus-also-accumulates-maintenance | Shared instructions require upkeep when different agent tools consume different files; repairing that divergence is part of maintaining the engineering environment. -->

The shared knowledge also requires maintenance. Instruction files apply through the monorepo\textquotesingle s directory hierarchy, while different agent tools read different files. Shopify describes divergence between those conventions: a directory can contain instructions for one tool and omit the corresponding file for another, so engineers working on the same code supply different context to their agents.\citep{lutke2026splitbrain, caswell2026agentsmd} Automation repairs that drift. Recording a lesson has therefore created another engineering artifact whose continued value depends on keeping it available and consistent. The cost follows from the mechanism that makes the knowledge reusable.

% <!-- point: shopify-converted-one-failure-into-a-code-enforced-guarantee-and-scopes-the-claim | A failed prompt-stated enumeration requirement motivated code enforcement in a documented security workflow; its admission arrangements and queue-use measurements do not characterize all River work. -->

A security-workflow account shows why recording an instruction may be insufficient. River violated a prompt-stated enumeration requirement and acknowledged the omission in its report. Shopify\textquotesingle s response was to require rules governing vulnerability closure to be enforced in code. The same account describes draft-only agent output, merge authority held by the engineers responsible for the affected code, a merge queue requiring fresh evidence, and a post-merge check that the default branch contains the fix. Use of that queue rose from about 10 to 80 percent.\citep{son2026security} That measures adoption of the queue, not the rate at which vulnerabilities escaped. The controls are described for security work; the account does not establish what evidence ordinary River pull requests must carry or how often they fail in production.

\section{Uber}\label{heading:uber}

% <!-- point: uber-treats-the-factory-itself-as-an-object-of-measurement-and-optimization | Uber treats foundation-model choice, context, orchestration, workflow cost, and reviewer attention as measurable engineering variables. -->

Uber treats the factory itself as something to measure and improve. Its managed agents perform work such as CI repair, review, alert triage, and maintenance; operational signals and schedules can initiate tasks as well as people.\citep{medisetty2026sdlc} The relevant question is what it costs to complete useful work through that arrangement. A token reduction, a cheaper foundation model, and a shorter review queue affect different parts of the answer.

% <!-- point: uber-tunes-the-environment-against-published-numbers | Context delivery, orchestration, foundation-model selection, and outcome-denominated costs make the agent's surrounding production system an object of optimization. -->

Uber adjusts context delivery, tool access, and coordination. Moving deterministic orchestration from repeated agent turns into executable code reportedly cut token use by more than half for common workflows and by more than 90 percent for bulk cases. Real workloads become benchmarks, foundation models are selected against cost and reliability, and managed agents are evaluated in units such as cost per merged pull request, review, alert, or cleanup.\citep{medisetty2026factory, medisetty2026sdlc} Those units connect resource use to an outcome rather than counting agent activity alone. They also make the surrounding engineering choices visible: some work is cheaper when executed directly than when repeatedly reconstructed through model interactions.

% <!-- point: uber-cannot-optimize-the-evidence-or-the-reviewer | Uber describes producer-assembled evidence and layered review with human admission; the public account does not establish independent re-derivation of the evidence or future automatic-approval criteria. -->

Evaluation begins inside the producing agent\textquotesingle s loop, with static-analysis repair, simulator screenshots compared with design specifications, and staging integration. The resulting pull request carries a table of checks and evidence assembled by that agent. The account does not say that independent machinery re-derives every entry.\citep{medisetty2026sdlc} Additional model-based review, CI, and human approval follow. Uber describes a possible move toward partial automatic approval, but does not publish the qualifying evidence or eligibility criteria.\citep{bond2026ureview} The reported review layers therefore tell us how evaluation is arranged without establishing the assurance of every proposed change.

% <!-- point: uber-publishes-cost-and-withholds-quality | Review latency and maintenance scheduling expose finite reviewer capacity; tracked quality measures and workflow costs do not by themselves establish required acceptance evidence. -->

Reviewer attention has its own capacity limit. Time to first review rose from three hours in 2024 to nine in 2026 as volume and change size grew.\citep{bond2026ureview} Uber schedules maintenance jobs on Sunday to use spare CI capacity, but separately caps the diffs engineers will face on Monday.\citep{medisetty2026sdlc} This is a concrete process-design decision: spare compute can produce more work than the next human stage can absorb. The organization also tracks revert rate, F1, and time to recovery alongside production cost, although the cited accounts do not publish a complete outcome series.

\section{Amazon}\label{heading:amazon}

% <!-- point: kiro-builds-engineering-artifacts-into-the-harness | Kiro supplies a reviewable specification workflow that adopting teams connect to their own repositories, tools, and engineering decisions. -->

Kiro supplies a reusable workflow for making intent explicit before implementation. A feature specification records user stories and acceptance criteria, develops an architectural design, and divides implementation into tasks that engineers can review.\citep{kiro2026specs} In MAGE\textquotesingle s vocabulary, these requirements and design artifacts are engineering models. DocAble arranged related work through project-specific work orders and design documents; Kiro makes a specification workflow part of the harness. Adopting it still leaves a team responsible for the content of those artifacts and their connection to its system.

% <!-- point: kiro-connects-requirements-to-tests-and-configurable-controls | Property-based tests provide evidence and configured hooks can enforce selected conditions; coverage and blocking decisions remain project responsibilities. | terms: modeling-principle, alignment-principle -->

The workflow can connect specifications to evidence and enforcement. Kiro\textquotesingle s IDE can extract properties from requirements and generate property-based tests, whose execution is optional by default and whose results are explicitly distinguished from proof.\citep{kiro2026correctness} Configured hooks can run checks after edits or block tool calls when preconditions fail.\citep{kiro2026hooks} Engineers must decide which properties matter, whether the tests express them correctly, and where failure should prevent work from proceeding. A supplied workflow makes these connections possible; its presence does not establish that a particular project enabled them.

% <!-- point: amazon-internal-use-adds-team-specific-production-structure | Bedrock's reported use combined Kiro with scoped work, tests, operational tools, and senior review; the twentyfold figure measures normalized commit velocity rather than quality or isolated causal effect. -->

Amazon\textquotesingle s Bedrock account makes some of that adopting-team work visible. Six senior engineers rebuilt part of the inference engine in 76 days, using small modules, reviewable commit-sized tasks, focused operational tools, tests, end-to-end canaries, and senior inspection of generated code.\citep{kiro2026bedrock} AWS reports a twentyfold increase in normalized commit velocity, from two to forty commits per developer per week.\citep{sivasubramanian2026frontier} That is a production-activity measure. It does not establish defect reduction or separate Kiro\textquotesingle s contribution from team selection and workflow redesign.

% <!-- point: amazon-governs-improvements-to-the-harness-itself | Kiro-trace tests proposed harness changes against an unchanged control and regression requirements, with human merge decisions and new checks derived from evaluation failures. | terms: tolerance, governance-conversion -->

Amazon also reports experiments on improving Kiro itself. Kiro-trace proposes prompt or tool changes from recorded runs. A separate deterministic gate screens targeted failures, checks regression assertions, and compares candidates with an unchanged control on the full benchmark. It requires gains on selected tasks without losses beyond a specified margin elsewhere; a person decides whether to merge. One apparent improvement was rejected after it proved to be a lucky run. Evaluation mistakes became checks, and confirmed fixes add regression assertions.\citep{kim2026kirotrace} These internal experiments illustrate the same engineering problem at another level: improving the fabricator requires evidence about the improvement and maintenance of the machinery that evaluates it.

\section{GitLab}\label{heading:gitlab}

% <!-- point: customer-intent-enters-as-files-and-customer-authority-as-a-withheld-capability | GitLab integrates organizational context with CI and permission boundaries, including a security-review flow that cannot approve changes. -->

GitLab illustrates how an integrated development workflow can bring existing organizational knowledge and controls to agent work. Repository instructions, chat rules, review skills, CI configuration, and workflow definitions supply different kinds of context.\citep{gitlab2026customize} The documented security-review flow can report findings but cannot approve a change.\citep{gitlab2026securityreview} For human-triggered flows running on runners, composite identity limits access to the intersection of the service account's and triggering user's privileges.\citep{gitlab2026identity} Tool-governance rules can also constrain actions; project rules may tighten but not loosen a group rule.\citep{gitlab2026toolgovernance} These are distinct functions: instructions inform the work, permissions restrict actions, and the approval boundary reserves a consequential decision.

GitLab\textquotesingle s internal playbook also makes acceptance depend on the surrounding pipeline. Its failing-test-first practice is backed by a CI job that fails when the test count decreases unless an explicit bypass label is present, addressing agents that remove tests to obtain a passing result.\citep{gitlab2026playbook} A test-count check cannot establish that the remaining tests express the right requirements. It does make one recurring way of evading them consequential at admission.

\section{What the Factories Share, and Where They Differ}\label{heading:what-the-factories-share-and-where-they-differ}

% <!-- point: the-factories-place-admission-in-four-different-locations | The cases make different admission arrangements visible and provide uneven evidence; reported activity should not be mistaken for independently established effectiveness. -->

The four factory questions reveal different arrangements rather than a common maturity sequence. The organizations differ in where they place engineering judgment, what they permit agents to decide, and which decisions remain with people or independent controls. Their measurements answer different questions: override frequency, veto and correction rates, merge rate, token cost, review delay, and engineering-health indicators cannot be combined into a single ranking of factory effectiveness.

% <!-- point: supervision-and-verification-do-not-scale-with-production | Production and supervisory capacity can grow at different rates; review delays, workload scheduling, and verification experience make this a process-design concern. -->

Two recurring costs cut across those arrangements. The first is supervisory capacity. Uber\textquotesingle s Sunday maintenance schedule would be incomplete without its separate limit on Monday\textquotesingle s review load. Spotify\textquotesingle s ability to generate more changes exposed limits in verification even when its measured rework did not rise. A factory therefore has to organize evaluation and admission alongside realization. Increasing the rate of one stage can shift delay or risk to another; the production system must be assessed through the work it can accept and sustain.

% <!-- point: the-inheritance-channel-carries-its-own-maintenance-cost | Shared knowledge and checks require maintenance; the documented revisions show how reusable engineering structure exchanges repeated judgment for upkeep. -->

The second cost is maintaining what later work inherits. Shopify repairs divergent instructions; Cloudflare revises policy representations and moves suitable checks earlier; Kiro\textquotesingle s improvement experiments acquire new regression assertions. Each preserved lesson can reduce future reconstruction or repeated judgment, but the resulting artifact must remain connected to the work it governs. This is the carrying-cost question from \magevalue{sec:the-software-factory} in operational form.
